%% GENERATED by python/paper/build.py from results/ -- do not edit by hand.
%% Rebuild:  .venv\Scripts\python.exe python\paper\build.py --only ArgentinaCalibration
%% Source:   results/paper/calibrationSummary.csv
\begin{table}[!htb]
\centering
\begin{threeparttable}
\caption{Calibration, Argentina}
\label{table:Arg:Calib}
\begin{tabular}{lll}
\toprule
\textbf{Parameter} & \textbf{Value} & \textbf{Target} \\
\midrule 
$\epsilon$ (\textit{pre}-reform) & $0.29$ & Minimum-to-full pension coverage \\[1ex]
$\theta$ & $0.84$ & Replacement rate dispersion \\[1ex]
$\alpha$ & $0.35$ & Factor income shares \\[1ex]
$\nu_t$ & $[1.65, 0.96]$ & 30-year gross population growth rates \\[1ex]
$\xi$ & $0.30$ & Elasticity of labor supply \\[1ex]
$\beta$ & $0.65$ & Capital--output ratio of $3.23$ \\[1ex]
$X$ & $1.94$ & Average formal workweek of 42.5 hours \\[1ex]
$\eta_i$ & $[0.64, 0.90, 1.14, 1.87]$ & Income distribution \\[1ex]
$\gamma_0$ & $0.32$ & Recipients of basic pension \\[1ex]
$\omega$ & $1.49$ & Social security tax of $10.9\%$ \\[1ex]
$\eta_0$ & $0.344$ & Informal relative income, eq.\ \eqref{eq:calibration_eta} \\[1ex]
$X_0$ & $0.804$ & Informal relative working hours \\[1ex]
\bottomrule
\end{tabular}
\begin{tablenotes}
\footnotesize
\item \textit{Note:} Our default specification relies on $\rho=1$. Relative formal hours are then a prediction: the model gives $[0.86, 0.95, 1.02, 1.18]$ against $[0.89, 0.98, 1.02, 1.10]$ in the data. Common-$X$ calibration: one leisure parameter $X$ shared across formal income groups, with the hours unit pinned by targeting the observed average workweek, so relative hours are a prediction rather than a calibration target. The informal $\eta_0$ and $X_0$ are calibrated under either variant.
\end{tablenotes}
\end{threeparttable}
\end{table}
